\section{训练数据构建：覆盖、组合爆炸与泛化}

\subsection{为什么数据问题在 Agent 里更难}
讲者把 ``高质量训练数据'' 放在第一优先级。原因在于 Agent 任务空间是组合爆炸的：环境类型、工具集合、验证方式、任务目标可自由组合。即便在单一领域（如 coding），也会出现仓库结构、测试风格、错误类型、依赖生态的大幅变化。

\begin{importantbox}{Agent 数据难点是组合，而非规模}
在 Agent 训练中，单纯增加样本数量并不能保证泛化。真正关键是是否覆盖了 ``环境-工具-验证器'' 的代表性组合，以及这些组合之间是否有足够多的跨域迁移信号。
\end{importantbox}

\subsection{数据集构建原则}
结合讲座内容，可以将可验证 Agent 的数据构建原则归纳为以下四点：任务类型广度、动作多样性、验证维度完整性、失败模式显式化。尤其是失败样本，往往比成功样本更能提升鲁棒性。

\begin{table}[H]
\centering
\begin{tabular}{p{2.4cm}p{3.7cm}p{5.4cm}}
\toprule
\textbf{原则} & \textbf{目标} & \textbf{落地做法} \\
\midrule
环境多样性 & 减少场景过拟合 & 同时引入数学、代码、检索、业务 API 等环境 \\
动作多样性 & 提升策略广度 & 对同一任务保留多条高质量轨迹而非单解 \\
验证多样性 & 防止单指标投机 & 单元测试、格式检查、规则引擎、人工抽检联合 \\
失败样本显式化 & 学习错误边界 & 记录失败原因并在训练中加入负反馈 \\
\bottomrule
\end{tabular}
\caption{可验证 Agent 数据构建的四条主线}
\end{table}

\subsection{从 ``看起来强'' 到 ``真实泛化''}
讲者反复提醒：如果训练集与评估集过于相似，模型会出现 ``指标高但真实任务崩''。这在 Agent 场景更严重，因为任务链路更长，任何一步分布偏差都会放大终局误差。

\begin{warningbox}{数据-评估同分布幻觉}
若训练和评估过于接近，会导致团队误判模型能力：
\begin{itemize}
    \item 在已见任务上持续提升，但在新任务上性能急剧下降；
    \item 模型学会针对 benchmark 的策略，而非学到一般能力；
    \item 部署后出现 ``离线好、在线差'' 的系统性偏差。
\end{itemize}
\end{warningbox}

\subsection{本章小结}
Agent 数据工程的重点是 ``覆盖正确的组合空间''。与其追求样本总量，不如优先保障环境、工具、验证器和失败模式的结构化多样性。

